\section{Express the Number}
\label{sec:b-express}
\source{ICPC 2020--21 Amritapuri Preliminary, problem 2 (PYQ)}{Easy--Medium}

\begin{reconbox}
Only the official solution slides survive (CodeDrills is offline). They say:
``we can use only odd powers; an even power can be written as a sum of odd
powers; if the top power $2^k>x$ has $k$ odd subtract it, if $k$ is even
subtract $2^{k-1}$ twice; once $n\le x$ write it as $A[1]$; impossible for odd
$n$ with $x=0$; complexity $O(\log n)$.'' The statement below is the
reconstruction that matches every line of the slides. The greedy was checked
against an exhaustive DP for all $n\le1500$ and nine values of $x$.
\end{reconbox}

\problemstatement{Given $n$ and $x$, write $n=A_1+A_2+\dots+A_m$ where
$0\le A_1\le x$ and every other term is an \emph{odd} power of two
($2,8,32,128,\dots$). Minimise $m$, or print $-1$ if impossible.}

\begin{patternbox}
Representations with a restricted set of powers of two: process the binary
expansion from the top bit, and rewrite forbidden powers with allowed ones
($2^{2j}=2^{2j-1}+2^{2j-1}$).
\end{patternbox}

\subsection*{Approach and intuition}

While $n>x$ we must spend terms. Let $2^k$ be the top bit of $n$.
\begin{itemize}
  \item $k$ odd: subtract $2^k$ (one term). Using smaller odd powers instead
        would need at least two terms for the same amount.
  \item $k$ even: $2^k$ is not allowed; the largest allowed power below $n$ is
        $2^{k-1}$. Subtract it (one term); after that the top bit is at most
        $k-1$ (odd) and the loop continues.
\end{itemize}
When $n\le x$, the rest becomes $A_1$. If $x=0$ and $n$ is odd we can never
finish, since all odd powers of two are even.

\subsection*{Algorithm}
\begin{algorithmic}[1]
\If{$x=0$ \textbf{and} $n$ odd} \Return $-1$ \EndIf
\State $m\gets0$
\While{$n>x$}
  \State $k\gets hb(n)$; $n\gets n-(k\text{ odd}\ ?\ 2^k:2^{k-1})$; $m\gets m+1$
\EndWhile
\State \Return $m+1$
\end{algorithmic}

\dryrun{$n=10$, $x=3$: top bit $8=2^3$ (odd), $n\to2\le3$: $m=1$, answer $2$
($10=2+8$). $n=10$, $x=0$: $10\to2\to0$, answer $3$ ($A_1=0$, $8$, $2$).
$n=7$, $x=0$: odd, $-1$. \checkmark}

\subsection*{C++17 implementation}
\cppfile{bitwise/09_express_the_number.cpp}

\begin{complexitybox}
\textbf{Time:} $O(\log n)$. \textbf{Space:} $O(1)$.
\end{complexitybox}

\begin{pitfallbox}
\begin{itemize}
  \item $2^0=1$ is an \emph{even} power; it is never allowed.
  \item Use \texttt{unsigned long long} and \texttt{1ULL << k} for large $n$.
\end{itemize}
\end{pitfallbox}
